\section{Terms as containers, and decision diagrams}\label{sec:containers}

A state is a closed Rocq term; a decision diagram stores sets of integer
vectors of one length. This section says how one becomes the other, and
why the way it is done is the theory of \emph{containers}~\cite{abbott05}:
the code is organised along that theory (\file{lib/skeleton/container.ml},
\file{lib/learn/learn.ml}, ``Serialisation''), which lets one encoding
serve systems whose states have several shapes.

\subsection{Containers}

A container $S \lhd P$ is a set of \emph{shapes} $S$ and, for each shape
$s$, a set of \emph{positions} $P(s)$. Its extension is the functor
\[
  [\![ S \lhd P ]\!](X) \;=\; \Sigma\, s \in S.\; (P(s) \to X),
\]
an element of which is a shape and a value in $X$ at each of its
positions. Strictly positive types are the initial algebras of
containers, and containers are closed under the usual constructions, the
free monad among them~\cite{abbott05}; with positions that carry types,
they are indexed containers~\cite{altenkirch15}.

\subsection{Three stages}

\paragraph{The term language} (\code{Container.view}). A closed term is a
constructor applied to parameters and fields, or something else (a value
that is not a constructor application). This stage knows nothing of the
rules.

\paragraph{Refinement by the rules} (\code{Container.split},
\code{Skeleton.classify}). A constructor is \emph{structural} if some rule
keeps it in place: it occurs in the rule's source and the same
constructor is at the same position of its target
(\cref{sec:sym-skeleton}). Write the state type as $T = \mu X.\, F_s(X) +
F_v(X)$, with $F_s$ the structural constructors and $F_v$ the others. Then
\[
  T \;\cong\; F_s^{*}(V),
\]
where $F_s^{*}$ is the free monad on $F_s$ and $V$ the terms whose head
is not structural. The free monad on a container is a container: its
shapes are the finite trees of structural constructors with holes, its
positions the holes. A term is therefore a shape (its top part made of
structural constructors), and a value at each position (the subterm
where that part stops), and the split is canonical: no value starts with
a structural constructor. Positions are \emph{paths} from the root (field
indices), not slot numbers.

The rules also classify each summand (a rule chain from a shape,
\code{Skeleton.kind}), and the three classes are the three kinds of map
on container elements:
\begin{itemize}
\item \textbf{local} (the shape stays): the identity on shapes, a new
  value at some positions;
\item \textbf{structural} (the shape changes, values are moved or
  dropped: $0 \mid 0 \to 0$, Proc's \code{do_par_end} and \code{do_assoc}):
  a container morphism, a map on shapes and a renumbering of positions;
\item \textbf{dynamic} (a computed value has structure, as a process that
  becomes a parallel composition): the value is itself an element of
  $F_s^{*}(V)$, and the state is flattened by the free monad's
  multiplication $F_s^{*}F_s^{*} \Rightarrow F_s^{*}$; the new shape is
  known only with the value.
\end{itemize}

\paragraph{Serialisation} (\file{lib/learn/learn.ml}). A shape becomes a
number, a position (a path) a vector index, a value a number (tables
built as the engine meets them, \code{0} meaning ``no position here''). A
state is the vector $[\textit{shape}; \textit{value at each position}]$,
padded to the system's width (8 positions to start with, doubled when a
state needs more). A path is numbered the first time a shape of the
system has it, in the order of its slots' use for the start shape
(\cref{sec:sym-learn}), so \emph{the same path has the same index in every
shape}: a step that changes the shape at one place leaves every value
elsewhere at its index, and the engine's relation copies those indices
instead of reading them.

\subsection{Decision diagrams are sets of container elements}

A list decision diagram (\cref{sec:sym-ldd}) over such vectors is a set
of serialised container elements, and its structure follows the
container's:
\begin{itemize}
\item its first level is the shape, a finite set: the diagram branches
  first on which shape a state has;
\item each further level is a position, and a node there holds the values
  that position takes; positions a shape does not have hold $0$, shared by
  all the states of that shape, so they cost one node;
\item the sharing of the diagram is the sharing of values among
  positions: a product of choices, one per position (independent
  processes), costs a sum of nodes, not a product;
\item a summand's relation is a relation on the positions it reads and
  writes, the others copied: on the container, a local map touches the
  positions at or below one place; a structural map, the shape and the
  positions it moves; a dynamic one, the shape and the positions of the
  structure the value became.
\end{itemize}
Because positions are paths, this locality is the container's: a change
below a place is a change of the vector below that place's indices. With
positions numbered as slots (left to right in the current shape, as until
10 November 2026), removing one leaf of the tree renumbered every leaf
after it, so a structural step rewrote values it did not touch, and the
engine had to learn it on whole states.

\subsection{Up to structural congruence}\label{sec:containers:upto}

\paragraph{Reorderings, found in the rules} (\code{Skeleton.reorderings}).
A rule is a \emph{reordering} if it has no premise, its label is
\code{None}, its source is a structure over pattern variables each used
once, and its target uses the same variables each once, otherwise
arranged: a container morphism that only permutes positions. That covers
commutation (Proc's \code{do_comm}), associativity (\code{do_assocl},
\code{do_assocr}), a swap or a rotation inside a structure. Only
reorderings whose inverse is one too are kept, so that they are
reversible; \code{do_par_end} ($0 \mid 0 \to 0$) is not one. When none is
found, or not the right ones, the command \code{MeBi Congruence lts
rules} names them by hand.

\paragraph{Proofs up to congruence: to design.} The engine can put every
reachable state in a class: the states it reaches by the congruence's
silent steps, and that reach it back. A certificate over pairs of classes
must answer the moves of \emph{every member} of a class, not only those of
a representative. Checking the representatives alone is unsound: with $b
\xrightarrow{\tau} a \xrightarrow{x} 0$ and $a \xrightarrow{\tau} b$, the
class $\{a, b\}$ would be ``matched'' by an inert state. This is weak
bisimulation up to a relation contained in mutual silent reachability,
with the whole class's moves checked, which makes it sound where ``up to
$\approx$'' is not~\cite{pous-sangiorgi11}. It covers weak goals only:
the silent path to the representative is not a one-step answer. The
certificate comes from the engine and is loaded as a proof term; Rocq
recomputes nothing (\cref{sec:optimisations}).
